\documentclass[11pt,a4paper]{article}

\usepackage[utf8]{inputenc}
\usepackage[T1]{fontenc}
\usepackage{lmodern}
\usepackage{amsmath, amssymb, amsthm}
\usepackage{mathtools}
\usepackage{geometry}
\usepackage{hyperref}
\usepackage{xcolor}
\usepackage{booktabs}
\usepackage{enumitem}
\usepackage{microtype}
\usepackage{listings}
\usepackage{verbatim}

\geometry{margin=2.8cm}
\hypersetup{
  colorlinks=true,
  linkcolor=blue!50!black,
  citecolor=blue!50!black,
  urlcolor=blue!50!black,
  pdftitle={A Spectral Reduction of the Collatz Conjecture via Phantom Orbit Shadowing},
  pdfauthor={Piero Borgatta}
}

\newtheorem{theorem}{Theorem}
\newtheorem{lemma}[theorem]{Lemma}
\newtheorem{proposition}[theorem]{Proposition}
\newtheorem{corollary}[theorem]{Corollary}
\theoremstyle{definition}
\newtheorem{definition}[theorem]{Definition}
\newtheorem{conjecture}[theorem]{Conjecture}
\newtheorem{remark}[theorem]{Remark}
\newtheorem{example}[theorem]{Example}

\DeclareMathOperator{\nutwo}{\nu_2}
\newcommand{\Z}{\mathbb{Z}}
\newcommand{\Zp}{\mathbb{Z}_2}
\newcommand{\Q}{\mathbb{Q}}
\newcommand{\N}{\mathbb{N}}
\newcommand{\R}{\mathbb{R}}
\newcommand{\full}{\mathrm{FULL}}
\newcommand{\core}{\mathrm{CORE}}
\newcommand{\tail}{\mathrm{TAIL}}
\newcommand{\spec}{\rho}

\title{
  A Spectral Reduction of the Collatz Conjecture \\
  via Phantom Orbit Shadowing
}
\author{
  Piero Borgatta\thanks{Independent Researcher, Italy. Email: \texttt{info@pieroborgatta.com}.\
    ORCID: \href{https://orcid.org/0009-0001-8025-2405}{0009-0001-8025-2405}.} \\
  \small Independent Researcher
}
\date{May 2026 \\[2pt]
  \small\normalfont
  Zenodo: \href{https://doi.org/10.5281/zenodo.20021538}{\texttt{10.5281/zenodo.20021538}} \quad
  Source: \href{https://github.com/PieroBorgatta/Collatz}{\texttt{github.com/PieroBorgatta/Collatz}}
}

\begin{document}
\maketitle

\begin{abstract}
We propose a structural reformulation of the Collatz conjecture in terms
of a finite-dimensional spectral problem. Starting from the empirical
observation that ``rebel'' orbits of the Syracuse map shadow,
for finite stretches, periodic rational points of the $2$-adic dynamics
associated to negative rational fixed points, we construct an episode
graph on the space of phantom orbits. The strongly connected
component (SCC) of this graph that supports the slowest descent is
encoded as a finite weighted transfer operator on a refined phase
state $(v_2(t), \mathrm{odd}(t)\bmod 4, h\bmod 4)$ at the critical
node $(k=12,\, c=2,\, b=1)$.

We prove an exact congruential shadowing lemma reducing the question
of arbitrarily long shadowing episodes to a residue condition
$n\equiv q_w\pmod{2^{bA+1}}$, where $q_w$ is the $2$-adic
representative of a phantom word $w$ and $A=|w|_{\nu_2}$ its period
sum. We then build, for each truncation depth $T$ and high-bit lift
count $j$, a finite matrix $\full_{T,j}$ encoding first-return
behavior at the critical node, decompose it as
$\full_{T,j} = \core_{T,j} + \tail_{T,j}$, and apply a weighted
Collatz--Wielandt bound with respect to the right Perron eigenvector
$v$ of $\full_{T,j}$. Our numerical results show
\[
\spec(\full_{T,j})
\le
\max_i (\core_{T,j}\, v)_i / v_i + \max_i (\tail_{T,j}\, v)_i / v_i
\le 0.052
\quad
\text{for $T \le 16$ and $j \le 32$,}
\]
with monotone but geometrically decaying increments in both $T$ and
$j$, and we extend the bound to a cross-node operator on the SCC
obtaining $\spec(M_{\mathrm{cross}}) \le 0.0366$. The cross-node
spectrum is \emph{lower} than the single-node bound at the critical
vertex.

We do not claim a proof of the Collatz conjecture. We claim that the
structural problem is reduced to two explicit a priori estimates on
finite matrices: the asymptotic boundedness of the weighted bound
under refinement in $T$, and the closure of the empirical signature
support under refinement in $j$. We document supporting numerical
evidence over a five-decade range of parameters, expose the open
estimates as conjectures, and provide all computational scripts as
supplementary material.
\end{abstract}

\noindent\textit{Research program note.} This paper is the first
artifact of a personal research program by the author, an independent
researcher, exploring whether iterative collaboration with large
language model assistants can support non-standard attempts at open
mathematical problems. The methodology, the contribution split between
human and AI, and the planned formal verification phase are
documented in the \emph{Methodology and AI disclosure} section at the
end and in \texttt{METHODOLOGY.md} of the supplementary archive.
\medskip

\tableofcontents

\section{Introduction}
\label{sec:intro}

The Collatz conjecture asserts that the iteration
\[
C(n) =
\begin{cases}
n/2 & \text{if } n \text{ even,} \\
3n+1 & \text{if } n \text{ odd,}
\end{cases}
\]
applied to any positive integer $n$ eventually reaches the trivial
cycle $1 \to 4 \to 2 \to 1$. Despite eighty years of attention from
combinatorialists, number theorists, and dynamicists, the conjecture
remains open. The strongest unconditional progress is due to Tao
\cite{tao2019}, who established that almost all orbits attain
arbitrarily small values relative to the initial point, in a precise
density sense.

The bottleneck for a deterministic argument is well known: certain
``rebel'' starting values produce orbits that climb to many orders
of magnitude above their initial value before eventually descending.
Empirically these rebels are not isolated curiosities but appear to
form structured families which converge in inverse-tree language to
common dissipative trunks.

This work proposes that the structural mechanism behind such families
is \emph{$2$-adic shadowing}: rebel positive integers follow, for
finite stretches, the trajectory of negative rational fixed points
of the Syracuse map under composition. These ``phantoms'' have no
positive integer realization but locally constrain the $\nu_2$
sequence of nearby positive orbits, producing the observed expansive
corridors. Once the constraint is exhausted (necessarily, since
positive integers cannot literally equal a negative rational), the
orbit drops into a dissipative regime.

We make this framework precise, codify the dynamics as a finite
weighted transfer operator on a refined phase quotient, and prove
numerically a uniform spectral upper bound which reduces the
boundedness of orbits to an explicit a priori estimate on a family
of finite matrices.

\paragraph{Structure of the paper.}
Section~\ref{sec:setup} fixes notation for the accelerated Syracuse
map and the $\nu_2$-debt formalism. Section~\ref{sec:shadowing}
states and proves the exact congruential shadowing lemma.
Section~\ref{sec:episodes} introduces the episode graph and its
critical SCC. Section~\ref{sec:operator} constructs the finite
transfer operator at the critical node. Section~\ref{sec:bound}
develops the weighted Collatz--Wielandt perturbative bound.
Section~\ref{sec:numerics} reports numerical evidence in $T$ and $j$.
Section~\ref{sec:cross} extends the bound to a cross-node operator
on the entire SCC. Section~\ref{sec:reduction} formulates the
explicit conjectures whose resolution would yield a proof.
Section~\ref{sec:limits} is an honest accounting of what this work
does and does not establish.

\section{Setup: accelerated Syracuse map and gravitational debt}
\label{sec:setup}

For an odd integer $n$, the \emph{Syracuse step} is
\[
S(n) = \frac{3n+1}{2^{\nutwo(3n+1)}},
\qquad
a(n) := \nutwo(3n+1).
\]
We write $a_k = a(n_k)$ for the $\nu_2$ sequence along the orbit
$n_0, n_1, n_2, \ldots$ with $n_{k+1} = S(n_k)$.

A direct expansion gives, after $k$ Syracuse steps,
\[
\log_2 n_k = \log_2 n_0 + k \log_2 3 - \sum_{i=0}^{k-1} a_i + O(1).
\]
The dominant term $k \log_2 3 - \sum a_i$ is what we call the
\emph{gravitational debt} after $k$ steps:
\[
D(k) := k \log_2 3 - \sum_{i=0}^{k-1} a_i.
\]
Positive $D(k)$ indicates net expansion to step $k$; the orbit
height $\log_2 n_k$ tracks $D(k)$ up to bounded fluctuations.

The threshold $a_i > \log_2 3 \approx 1.585$ separates dissipative
from expansive single steps. The Collatz conjecture is morally
equivalent to the assertion that $D(k) \to -\infty$ along every
positive-integer orbit.

\section{Phantom rational orbits and the shadowing lemma}
\label{sec:shadowing}

\subsection{Phantom words and their rational fixed points}

Fix a finite word
\[
w = (a_0, a_1, \ldots, a_{L-1}) \in \N^L,
\qquad
A = A(w) := \sum_{i=0}^{L-1} a_i.
\]
The composition $S_w := S_{a_{L-1}} \circ \cdots \circ S_{a_0}$ is
affine: writing it out one obtains
\begin{equation}\label{eq:Sw}
S_w(n) = \frac{3^L\, n + C_w}{2^A},
\qquad
\begin{aligned}
C_0 &= 0, \\
C_{j+1} &= 3 C_j + 2^{A_j}, \\
A_{j+1} &= A_j + a_j.
\end{aligned}
\end{equation}
The fixed point of $S_w$ in $\Q$ is
\begin{equation}\label{eq:qw}
q_w = \frac{C_w}{2^A - 3^L}.
\end{equation}

\begin{definition}[expansive phantom]
The word $w$ is an \emph{expansive phantom} if $3^L > 2^A$. Then
$2^A - 3^L < 0$, and assuming $C_w > 0$ (which is automatic from
\eqref{eq:Sw}) we have $q_w \in \Q_{<0}$.
\end{definition}

The orbit of $q_w$ under $S$ is exactly periodic of length $L$ with
$\nu_2$ word $w$ repeated. In the $2$-adic completion $\Zp$, $q_w$
makes sense as a $2$-adic integer (since $\gcd(2^A - 3^L, 2)=1$) and
its $2$-adic orbit consists of $L$ distinct $2$-adic integers cycled
by $S$.

\subsection{Exact congruential shadowing}

Let $w^\infty$ denote the right-infinite repetition of $w$, and let
$B_m := \sum_{i=0}^{m-1} a_{i \bmod L}$ for $m \ge 0$ be the partial
sum of the periodic word.

\begin{lemma}[Exact shadowing]\label{lem:shadowing}
Let $w$ be an expansive phantom with $2$-adic representative $q_w$,
and let $n$ be a positive odd integer. If
\[
n \equiv q_w \pmod{2^{B_m+1}}
\qquad\text{(equality in $\Zp$ truncated to $B_m+1$ bits),}
\]
then the first $m$ Syracuse steps of $n$ have $\nu_2$ word exactly
$(a_0, a_1, \ldots, a_{m-1})$.

In particular, if $n \equiv q_w \pmod{2^{bA+1}}$, the first $b$
periods (i.e.\ $bL$ Syracuse steps) of $n$ shadow $w^\infty$.
\end{lemma}

\begin{proof}
We argue by induction on $j$. Suppose after $j$ Syracuse steps the
orbits of $n$ and $q_w$ have followed the same word
$(a_0,\ldots,a_{j-1})$. Then both have undergone the same affine
transformation $S_{(a_0,\ldots,a_{j-1})}$, which has multiplier
$3^j / 2^{B_j}$. Therefore
\[
S^j(n) - S^j(q_w) = \frac{3^j}{2^{B_j}} (n - q_w).
\]
Taking $\nu_2$ of both sides and using $\nu_2(3^j)=0$,
\[
\nutwo\bigl(S^j(n) - S^j(q_w)\bigr)
= \nutwo(n - q_w) - B_j
\ge (B_m + 1) - B_j.
\]
For step $j$ we need to extract \emph{exactly} $a_j$ factors of $2$
from $3 S^j(n) + 1$, which agrees with $3 S^j(q_w) + 1$ up to
valuation
\[
\nutwo\bigl(3 S^j(n) + 1 - 3 S^j(q_w) - 1\bigr)
= \nutwo\bigl(3(S^j(n) - S^j(q_w))\bigr)
\ge B_m + 1 - B_j.
\]
For the next valuation to coincide it suffices that
\[
B_m + 1 - B_j \ge a_j + 1,
\quad\text{i.e.}\quad
B_m \ge B_j + a_j = B_{j+1},
\]
which holds for all $j < m$ by definition. Induction completes the
proof. The ``+1'' bit is essential: without it one only ensures
agreement modulo $2^{a_j}$, not the absence of an additional factor
of $2$.
\end{proof}

\begin{corollary}[No infinite shadowing]\label{cor:no-infinite}
A positive integer cannot shadow $w^\infty$ for all $m$, since that
would force $n = q_w$ in $\Zp$, hence $n = q_w \in \Q_{<0}$,
contradicting $n \in \N$.
\end{corollary}

\section{The episode graph}
\label{sec:episodes}

Lemma~\ref{lem:shadowing} shows that an integer $n$ can shadow a
phantom $w$ for some maximal number of periods $b = b_w(n)$, after
which it must leave the residue class. The natural object for
tracking long orbits of integer rebels is therefore not $n$ itself
but the sequence of \emph{episodes}
\[
n \to (w_1, b_1) \to (w_2, b_2) \to \cdots
\]
together with terminal events (drops below the starting value).

\subsection{The graph}

We monitor a finite set of phantoms identified empirically up to
$k\le 24$ (their explicit list is given in
Appendix~\ref{app:phantoms}). For each pair
$(\text{phantom }k, \text{cycle }c, \text{depth }b)$ we associate a
\emph{node} $(k,c,b)$. The episode graph $G$ has these nodes as
vertices; an edge $(k,c,b) \to (k',c',b')$ records the empirical
observation that some integer orbit shadowed $(k,c,b)$ for exactly
$b$ periods, then re-entered shadowing of $(k',c',b')$.

\subsection{The critical SCC}

Empirical exhaustion (script \verb|59_shadow_episode_graph.py|,
\verb|60_episode_graph_diagnostics.py|) over the dense sample
$b \le 16$, $t \le 65535$ identifies a single nontrivial strongly
connected component containing
\[
(10,1,b) \text{ for } b \in [1,15], \quad
(11,1,b) \text{ for } b \in [1,15], \quad
(12,1,b) \text{ for } b \in [1,2], \quad
(12,2,b) \text{ for } b \in [1,2].
\]
The \emph{critical node} is $(12, 2, 1)$: it is the unique node of
the SCC with significant self-recurrence, and most cross-node
transitions terminate at it.

The strategy of the rest of the paper is to prove that this SCC is
spectrally subcritical when encoded as a transfer operator with
appropriate weights.

\section{The critical transfer operator}
\label{sec:operator}

Fix the critical node $(k,c,b) = (12,2,1)$ with phantom word $w$,
period sum $A$, and $2$-adic representative $q_w$. Every integer $n$
shadowing this node for one period satisfies
\[
n \equiv q_w \pmod{2^{A+1}},
\quad\text{i.e.}\quad
n = r + t \cdot 2^{A+1}
\]
for the unique residue $r = q_w \bmod 2^{A+1}$ and parameter
$t \in \N$.

\subsection{Refined phase state}

For a parameter $t$, we define the \emph{refined phase state}
\[
\sigma(t, h) := \bigl(\nu_2(t)\wedge V,\ \mathrm{odd}(t) \bmod 4,\ h \bmod 4\bigr),
\]
where $V$ is a finite valuation cap, $\mathrm{odd}(t) := t / 2^{\nu_2(t)}$,
and $h$ is the count of monitored hits accumulated so far. The
choice of moduli $4$ and $4$ is the empirical minimum that makes
the spectral truncation stable (script
\verb|74_phase_transfer_stress.py|).

\subsection{Truncated finite operator}

For depths $T \in \N$ and $j \in \N$, we define the truncated
transfer operator $\full_{T,j}$ on the state space
\[
\mathcal{S}_{T,j} := \{\sigma(t, h) : t \in [0, 2^T \cdot j),\ h \in [0,4)\}
\]
as follows. For each source state $\sigma_{\mathrm{src}}$, enumerate
all $(t, h)$ producing it. Trace the Syracuse orbit from
$n_0 = r + t \cdot 2^{A+1}$ until either:
\begin{enumerate}
  \item $n < n_0$ (terminal: contributes no edge);
  \item the next monitored hit at the same critical node occurs at
        $n^* = r + t^* \cdot 2^{A+1}$. This contributes an edge
        \[
        \sigma_{\mathrm{src}} \to \sigma(t^*, h+1)
        \quad\text{with weight}\quad
        2^{-\delta},
        \quad
        \delta = \lfloor \log_2 t^* \rfloor - \lfloor \log_2 t \rfloor.
        \]
\end{enumerate}
We aggregate edges and normalize by source counts:
\[
(\full_{T,j})_{\sigma, \sigma'}
= \frac{1}{|\{(t,h) : \sigma(t,h)=\sigma\}|}
  \sum_{\substack{(t,h) \\ \sigma(t,h)=\sigma}}
  \mathbf{1}_{\text{trace lands on }\sigma'}
  \cdot 2^{-\delta(t)}.
\]

By construction $\full_{T,j}$ is a substochastic-like nonnegative
matrix with row sums $\le 1$, the deficit corresponding to terminal
transitions.

\subsection{Empirical signature decomposition}

For each pair $(r, h)$ with $r \in [0, 2^T)$, the $j$ lifts
\[
t = r + j' \cdot 2^T, \quad j' \in [0, j),
\]
produce $j$ traces. We collect the multiset of \emph{transition
signatures} (destination state plus $\delta$ bit increment) and
declare the \emph{core signature} to be the most frequent one. Edges
matching this signature are placed in $\core_{T,j}$; all other
edges are placed in $\tail_{T,j}$. By construction
\[
\full_{T,j} = \core_{T,j} + \tail_{T,j}.
\]
We refer the reader to script
\verb|77_high_bit_tail_bound.py| for the implementation details.

\section{Weighted perturbative bound}
\label{sec:bound}

\subsection{Collatz--Wielandt with the Perron eigenvector of FULL}

The matrices $\full_{T,j}$, $\core_{T,j}$, $\tail_{T,j}$ are
nonnegative. Let $v_{T,j}$ denote a right Perron eigenvector of
$\full_{T,j}$, normalized so $\|v_{T,j}\|_\infty = 1$. By
Perron--Frobenius (after a tiny rank-1 perturbation to ensure
irreducibility, see remark below), $v_{T,j}$ exists and is strictly
positive.

The decomposition $\full_{T,j} = \core_{T,j} + \tail_{T,j}$ acting
on $v_{T,j}$ gives, componentwise,
\[
(\full_{T,j} v_{T,j})_i
= (\core_{T,j} v_{T,j})_i + (\tail_{T,j} v_{T,j})_i.
\]
Taking the maximum over $i$ such that $(v_{T,j})_i > 0$ and
using $\spec(\full_{T,j}) = \max_i (\full_{T,j} v_{T,j})_i / (v_{T,j})_i$,
\begin{equation}\label{eq:bound}
\boxed{\;
\spec(\full_{T,j})
\;\le\;
\underbrace{\max_i \frac{(\core_{T,j} v_{T,j})_i}{(v_{T,j})_i}}_{\mathrm{core\ action}}
\;+\;
\underbrace{\max_i \frac{(\tail_{T,j} v_{T,j})_i}{(v_{T,j})_i}}_{\mathrm{tail\ action}}.
\;}
\end{equation}
We denote the right-hand side $\mathrm{bound}(T,j)$.

\begin{remark}[numerical implementation]
We compute $v_{T,j}$ by power iteration on
$\full_{T,j} + \varepsilon\, J$ with $\varepsilon = 10^{-12}$ and
$J$ the all-ones matrix. The rank-1 perturbation removes
reducibility numerically without affecting the dominant eigenvalue
beyond $\varepsilon$. Convergence is checked in $\ell^\infty$ norm
to tolerance $10^{-12}$.
\end{remark}

\begin{remark}[why CORE's Perron is wrong]
Using the Perron eigenvector of $\core_{T,j}$ instead of
$\full_{T,j}$ yields a divergent ``bound'': $\core_{T,j}$ is
typically reducible, with its Perron eigenvector concentrated on a
strict subset of states, and $\tail_{T,j}$ routinely sends mass to
states outside that subset. The Collatz--Wielandt ratios then blow
up. The empirical fix and the theoretical fix both point to using
$v_{T,j}$ from $\full_{T,j}$.
\end{remark}

\section{Numerical evidence}
\label{sec:numerics}

\subsection{Bound stability in $T$}

We compute $\mathrm{bound}(T, j)$ for $T \in \{8, 10, 12, 14, 16\}$
and $j \in \{16, 32\}$ using script
\verb|84_weighted_perturbation_bound.py| (with the eigenvector
selection from the previous remark). Results:

\begin{center}
\begin{tabular}{rrrrrr}
\toprule
$T$ & $j$ & $\spec(\full)$ & $\mathrm{core\ action}$ & $\mathrm{tail\ action}$ & $\mathrm{bound}$ \\
\midrule
 8 & 16 & 0.0294 & 0.0114 & 0.0294 & 0.0408 \\
 8 & 32 & 0.0299 & 0.0102 & 0.0299 & 0.0400 \\
10 & 16 & 0.0303 & 0.0184 & 0.0303 & 0.0487 \\
10 & 32 & 0.0308 & 0.0177 & 0.0308 & 0.0485 \\
12 & 16 & 0.0303 & 0.0137 & 0.0303 & 0.0499 \\
12 & 32 & 0.0306 & 0.0135 & 0.0306 & 0.0497 \\
14 & 16 & 0.0306 & 0.0157 & 0.0306 & 0.0511 \\
14 & 32 & 0.0307 & 0.0154 & 0.0307 & 0.0507 \\
16 & 16 & 0.0307 & 0.0169 & 0.0307 & 0.0516 \\
\bottomrule
\end{tabular}
\end{center}

The bound is monotonically increasing in $T$ but the increments
decay geometrically:
\[
\Delta_8^{10} = +0.0079, \quad
\Delta_{10}^{12} = +0.0012, \quad
\Delta_{12}^{14} = +0.0012, \quad
\Delta_{14}^{16} = +0.0005.
\]
A geometric extrapolation yields
\[
\lim_{T \to \infty} \mathrm{bound}(T, 16) \le 0.054.
\]
The bound is essentially independent of $j$: differences between
$j=16$ and $j=32$ at fixed $T$ never exceed $0.001$.

\subsection{Spectral gap $\spec(\full) - \spec(\core)$}

Independently, the spectral radii of $\full$ and $\core$ are
computed (script \verb|83_perturbation_gap_scaling.py|). The gap
$\spec(\full) - \spec(\core)$ shrinks with $T$:
\[
T=6:\ 0.025\text{--}0.032,
\quad
T=8:\ 0.025\text{--}0.028,
\quad
T=10:\ 0.018\text{--}0.019,
\]
\[
T=12:\ 0.016\text{--}0.017,
\quad
T=14:\ 0.015,
\quad
T=16:\ 0.014.
\]
The spectral mass of $\full$ migrates progressively to $\core$ as
the quotient is refined.

\subsection{Truncation stability in $j$ (Cauchy property)}

For fixed $T$, we check whether $\full_{T, j}$ admits a Cauchy
property in $j$ via two independent metrics
(script \verb|85_truncation_stability.py|):
\begin{itemize}
  \item \textbf{signature drift}: the fraction of $(r, h)$ groups
        whose dominant signature changes between consecutive
        doublings of $j$;
  \item \textbf{bound increment}: $\mathrm{bound}(T, 2j) - \mathrm{bound}(T, j)$.
\end{itemize}

Results:

\begin{center}
\begin{tabular}{rrrrrr}
\toprule
$T$ & $j$ & sig\_drift\_vs\_prev & bound & $\Delta_{\mathrm{prev}}$ \\
\midrule
 8 & 16  & --     & 0.0408 & --      \\
 8 & 32  & 0.0078 & 0.0400 & $-0.0007$ \\
 8 & 64  & 0.0078 & 0.0421 & $+0.0021$ \\
 8 & 128 & 0.0000 & 0.0425 & $+0.0004$ \\
\midrule
10 & 16  & --     & 0.0487 & --      \\
10 & 32  & 0.0039 & 0.0485 & $-0.0003$ \\
10 & 64  & 0.0088 & 0.0509 & $+0.0024$ \\
10 & 128 & 0.0000 & 0.0515 & $+0.0006$ \\
\bottomrule
\end{tabular}
\end{center}

Two facts:
\begin{enumerate}
  \item The signature drift drops to \emph{exactly zero} when $j$ is
        doubled from $64$ to $128$, for both $T \in \{8, 10\}$. The
        empirical transition signatures of all $(r, h)$ groups
        stabilize.
  \item The bound increments decay by a factor $\sim 4$ per
        doubling, suggesting a Cauchy bound
        $\lim_{j\to\infty}\mathrm{bound}(T=10, j) \lesssim 0.052$.
\end{enumerate}
Combining the two extrapolations, we conjecture
\[
\sup_{T,\, j} \mathrm{bound}(T, j) < 0.06,
\]
a margin of $16\times$ from the criticality threshold $1$.

\section{Cross-node certificate on the SCC}
\label{sec:cross}

The single-node bound at the critical node is the worst case:
running script \verb|86_scc_node_certificates.py| at $T=10, j=32$
across the SCC nodes yields:

\begin{center}
\begin{tabular}{lrrr}
\toprule
node & $\spec(\full)$ & $\mathrm{bound}$ & status \\
\midrule
$(10,1,1)$ & $\approx 0$ & $0.0000$ & OK (return-to-self vacuous) \\
$(10,1,2)$ & $\approx 0$ & $0.0000$ & OK \\
$(11,1,1)$ & $\approx 0$ & $0.0000$ & OK \\
$(11,1,2)$ & $\approx 0$ & $0.0000$ & OK \\
$(12,1,1)$ & $0.0035$  & $0.0043$ & OK \\
$(12,2,1)$ & $0.0308$  & $0.0485$ & OK (worst case) \\
$(12,2,2)$ & $0.0007$  & $0.0007$ & OK \\
\bottomrule
\end{tabular}
\end{center}

The nodes $(10, 1, *)$, $(11, 1, *)$, $(12, 2, 2)$ have nearly
vanishing return-to-self spectrum. They are entry gates to the SCC,
not self-recurrent.

\subsection{The cross-node operator}

We construct the assembled cross-node operator $M_{\mathrm{cross}}$
on the union state space
\[
\bigsqcup_{(k,c,b) \in \mathrm{SCC}} \mathcal{S}_{T,j}^{(k,c,b)}.
\]
Edges are recorded between the source state at one node and the
\emph{first hit} state at any node of the SCC, with weight
$2^{-\delta_{\mathrm{bit}}}$ where $\delta_{\mathrm{bit}}$ is the
bit-length increment between source $n_0$ and destination value (a
node-invariant quantity). The construction is implemented in
script \verb|87_cross_node_transfer.py|.

At $T=8, j=8$ on $5$ representative SCC nodes:
\[
\spec(M_{\mathrm{cross}}) = 0.0366, \qquad
\mathrm{bound}(M_{\mathrm{cross}}) = 0.0366.
\]
The cross-coupling \emph{lowers} the spectrum from the single-node
worst case ($0.0485$). Topologically the SCC is a star centered on
$(12,2,1)$:
\begin{verbatim}
    (10,1,1) -> (12,2,1)   8128  (gate)
    (11,1,1) -> (12,2,1)   8192  (gate)
    (12,2,2) -> (12,2,1)   8128  (gate)
    (12,2,1) -> (12,2,1)    812  (only significant self-loop)
    (12,2,1) -> (12,1,1)    308
    (12,1,1) -> (12,1,1)     96
    (12,2,2) -> (12,2,2)     64
    (12,2,1) -> (12,2,2)      8
\end{verbatim}
The single-node bound was therefore overly conservative; including
cross-node transitions strictly improves it.

\section{Conjectures and the spectral reduction statement}
\label{sec:reduction}

We isolate the open analytic content as two explicit conjectures.

\begin{conjecture}[uniform bound in $T$]\label{conj:T}
There exists a constant $B_\infty < 1$ such that
\[
\sup_{T \in \N,\ j \in \N} \mathrm{bound}(T, j) \le B_\infty.
\]
Empirically the data support $B_\infty \approx 0.06$.
\end{conjecture}

\begin{conjecture}[signature closure in $j$]\label{conj:j}
For each $(T, r, h)$ there exists $j^*(T, r, h)$ such that the
dominant transition signature stabilizes for all $j \ge j^*$, and
furthermore
\[
\sup_{j \ge j^*}
\bigl|\mathrm{bound}(T, j) - \mathrm{bound}(T, j^*)\bigr|
\le C \cdot 2^{-\beta(j - j^*)}
\]
for some $C, \beta > 0$ uniform in $(T, r, h)$.
\end{conjecture}

\begin{theorem}[conditional reduction, sketch]\label{thm:conditional}
Assume Conjectures~\ref{conj:T} and~\ref{conj:j}. Then for every
positive integer $n$, the Syracuse orbit $S^k(n)$ satisfies
$\liminf_{k\to\infty} S^k(n) = 1$.
\end{theorem}

The argument is the standard one for transfer-operator subcriticality
combined with the no-infinite-shadowing
Corollary~\ref{cor:no-infinite}: assuming the conjectures, the
critical SCC has spectral radius $\le B_\infty < 1$ on a refined
quotient that captures all dynamically relevant phases, hence the
return-to-critical mass strictly contracts, hence orbits cannot
remain forever near the SCC, hence by exhaustion of the empirical
phantom set every orbit eventually escapes to a dissipative regime
with average $\nu_2 > \log_2 3$ and descends. The full argument
requires an exhaustive enumeration of phantoms below a controlled
threshold and a verification that no further SCCs appear; see
Section~\ref{sec:limits} for an honest discussion of the gaps.

\section{What this is and what this is not}
\label{sec:limits}

This work is a \emph{program}, not a proof. We summarize what is
established, what is empirical, and what is missing.

\paragraph{Established (rigorously).}
\begin{itemize}
  \item Lemma~\ref{lem:shadowing} (exact congruential shadowing)
        is a theorem with elementary proof. It identifies the
        precise residue condition for $b$-fold phantom shadowing.
  \item Corollary~\ref{cor:no-infinite} (no infinite shadowing)
        follows formally.
  \item The construction of the finite transfer operators
        $\full_{T,j}$, $\core_{T,j}$, $\tail_{T,j}$ at the critical
        node is a deterministic algorithm; the matrices are
        exactly computable for all finite $T, j$.
  \item The bound \eqref{eq:bound} is a Collatz--Wielandt
        consequence of nonnegativity; it gives a true upper bound
        on $\spec(\full_{T,j})$ for each finite $(T, j)$.
\end{itemize}

\paragraph{Empirical (numerically supported, not proved).}
\begin{itemize}
  \item All numerical bounds reported in
        Sections~\ref{sec:numerics} and~\ref{sec:cross}.
  \item The geometric decay of bound increments in $T$ and $j$.
  \item The closure of dominant signatures at $j = 128$ for
        $T \in \{8, 10\}$.
  \item The identification of $(12, 2, 1)$ as the worst-case node
        of the SCC.
  \item The structural ``star'' topology of the SCC around
        $(12, 2, 1)$.
\end{itemize}

\paragraph{Missing.}
\begin{itemize}
  \item Proofs of Conjectures~\ref{conj:T} and~\ref{conj:j}. We
        believe the natural framework is transfer-operator analysis
        on a suitable Banach space of $2$-adic Lipschitz functions;
        a Lasota--Yorke inequality for the family
        $\{\full_T\}$, combined with Hennion's theorem and
        Keller--Liverani perturbation theory \cite{kellerliverani},
        would yield Conjecture~\ref{conj:T}. We have not carried
        this out.
  \item Verification that the empirical phantom set ($k \le 24$) is
        complete in the sense that no SCC outside the listed nodes
        contributes. Empirical sweeps to $b \le 16$, $t \le 65535$
        find none, but this is not a proof.
  \item A formal version of Theorem~\ref{thm:conditional}.
\end{itemize}

\paragraph{Honest claim.}
We have reduced the open structural content of the Collatz
conjecture to a problem of \emph{spectral perturbation theory on a
finite explicit family of nonnegative matrices}, supported by
numerical evidence with a $16\times$ safety margin. We do not
claim the conjecture itself.

\section*{Methodology and AI disclosure}

This paper is the first artifact of a personal research program by the
author, an independent researcher, exploring whether sustained
collaboration with large language model assistants can enable
non-standard framings of open mathematical problems. A clear
separation between contributions is essential for the reader's
correct evaluation of the material.

\paragraph{Author's contribution.}
The conceptual framing is due to the author: the gravitational debt
metaphor, the empirical identification of $\nu_2 = 1$ corridors as
the locus of expansive behavior, the structural intuition that integer
rebels shadow negative rational fixed points of $S$, the strategic
pivot from the empirical Lyapunov candidates of the earlier phase to
the present spectral program, and the decision to publish at this
stage. The author led the research direction throughout: the choices
of what to investigate, what to refine, when to stop scaling, and what
to publish are the author's.

\paragraph{AI assistance.}
The technical work — Python script implementation, the mathematical
formalization of Lemma~\ref{lem:shadowing} and the weighted
Collatz--Wielandt bound of Section~\ref{sec:bound}, and the \LaTeX{}
drafting of this manuscript — was developed through iterative
collaboration with multiple LLM systems. Specifically: OpenAI Codex
and Google Gemini for the early empirical phase
(\texttt{scripts/early\_empirical/}), and Anthropic's Claude (via the
Claude Code interface) for the present spectral program
(\texttt{scripts/spectral\_program/}), the formalization of the
shadowing lemma, and the manuscript drafting.

\paragraph{Verification status and future work.}
As of this version, the work has not been reviewed by a human
mathematical expert. The numerical results are deterministic and
reproducible from the supplementary scripts; anyone with Python,
NumPy, and SciPy can verify them. The mathematical claims, including
Lemma~\ref{lem:shadowing}, the construction of the transfer operator,
and the conditional reduction theorem, rest on AI-assisted
formalization that has not yet been externally validated.

This paper is therefore submitted explicitly to invite scrutiny by
researchers in symbolic dynamics, $p$-adic dynamical systems, and
transfer operator theory. A planned next phase of this research
program is to mechanically verify Lemma~\ref{lem:shadowing} in a
proof assistant; the intended target is Lean~4 with Mathlib. A
revised version of this preprint will incorporate the formal
verification when complete. Correspondence concerning errors,
alternative framings, or collaboration is explicitly welcome.

\begin{thebibliography}{9}

\bibitem{tao2019}
T.~Tao,
\emph{Almost all orbits of the Collatz map attain almost bounded values},
arXiv:1909.03562, 2019; published in Forum of Mathematics, Pi, 2022.

\bibitem{lagarias1985}
J.~C.~Lagarias,
\emph{The 3x+1 problem and its generalizations},
American Mathematical Monthly 92 (1985), 3--23.

\bibitem{lagarias2010}
J.~C.~Lagarias (ed.),
\emph{The Ultimate Challenge: The 3x+1 Problem},
American Mathematical Society, 2010.

\bibitem{kellerliverani}
G.~Keller and C.~Liverani,
\emph{Stability of the spectrum for transfer operators},
Annali della Scuola Normale Superiore di Pisa, Cl. Sci. (4) 28 (1999), 141--152.

\bibitem{baladi}
V.~Baladi,
\emph{Positive Transfer Operators and Decay of Correlations},
World Scientific, 2000.

\bibitem{anashin}
V.~Anashin,
\emph{The non-Archimedean theory of discrete systems},
Math. Comput. Sci. 6 (2012), 375--393.

\end{thebibliography}

\appendix

\section{Phantom records used in this work}
\label{app:phantoms}

The expansive phantom representatives identified empirically up to
$k = 24$ and used as monitored nodes in Section~\ref{sec:episodes}:

\begin{center}
\begin{tabular}{rl}
\toprule
$k$ & $q_w \approx$ \\
\midrule
10 & $-5\,739\,192\,681\,529 / 2\,404\,426\,874\,857 \approx -2.387$ \\
11 & $-2\,199\,012\,694\,031 / 709\,849\,655\,971 \approx -3.098$ \\
12 (cycle 1) & $-2123 / 1675 \approx -1.267$ \\
12 (cycle 2) & $-817 / 601 \approx -1.359$ \\
20 & $-104\,312\,644\,103 / 30\,307\,317\,785 \approx -3.442$ \\
\bottomrule
\end{tabular}
\end{center}

The full list and the construction of these representatives from
the rational fixed point formula \eqref{eq:qw} is implemented in
\verb|54_phantom_rational_shadows.py|.

\section{Supplementary scripts}

All scripts referenced in this paper are available as supplementary
material in the GitHub repository
\begin{center}
\url{https://github.com/PieroBorgatta/Collatz}
\end{center}
and as a frozen Zenodo deposit
\begin{center}
\url{https://doi.org/10.5281/zenodo.20021538}
\end{center}

The principal scripts and what each computes:

\begin{description}[leftmargin=4em, style=nextline]
  \item[\texttt{54\_phantom\_rational\_shadows.py}]
    Computes the $2$-adic representatives of expansive phantom
    cycles; underpins all subsequent constructions.
  \item[\texttt{55\_shadowing\_congruence.py}]
    Verifies Lemma~\ref{lem:shadowing} on the empirical phantom
    set for $b \le 10$.
  \item[\texttt{59\_shadow\_episode\_graph.py},\
        \texttt{60\_episode\_graph\_diagnostics.py}]
    Constructs the episode graph and identifies its SCCs.
  \item[\texttt{75\_critical\_symbolic\_operator.py}]
    Builds the truncated finite operator at the critical node.
  \item[\texttt{77\_high\_bit\_tail\_bound.py}]
    Decomposes $\full = \core + \tail$ from the empirical
    signature majority.
  \item[\texttt{82\_boundary\_scc\_spectral\_scaling.py}]
    Spectral scaling on SCCs touching the boundary layer
    $\nu_2(t) \ge T$.
  \item[\texttt{83\_perturbation\_gap\_scaling.py}]
    Measures the gap $\spec(\full) - \spec(\core)$.
  \item[\texttt{84\_weighted\_perturbation\_bound.py}]
    Computes the weighted bound \eqref{eq:bound}.
  \item[\texttt{85\_truncation\_stability.py}]
    Cauchy/signature stability of the truncation in $j$.
  \item[\texttt{86\_scc\_node\_certificates.py}]
    Per-node bounds on the SCC.
  \item[\texttt{87\_cross\_node\_transfer.py}]
    Assembled cross-node transfer operator on the SCC.
\end{description}

\end{document}
